\section{相关工作}
\label{sec:related}

\subsection{多模态问答与主动澄清}

视觉问答中的歧义与澄清是近年研究热点。Jian 等人提出 ClearVQA 基准与自动生成"歧义--澄清"问题对的管线，解决用户提问中"指代不明、含义不清"的歧义，并训练 VLM 先澄清再作答~\cite{jian2025teaching}；PromptCap 则将澄清建模为任务生成问题~\cite{hu2023promptcap}。这些工作共同表明 VLM 天生"爱答不爱问"。\textbf{本文的区别}在于：它们解决的是"问题歧义"，而 AquaDiag 解决的是"诊断信息不足"——模型需要追问的是缺失的临床字段（水温、死亡率、粪便），而非澄清指代对象；此外，我们额外叠加了\textbf{合规红线}与\textbf{拒答判断}两个维度。

\subsection{医疗/诊断对话与槽位填充}

医疗对话系统长期以槽位填充方式组织问诊流程。MedDG 构建了多领域中文医疗对话数据集与槽位填充评测~\cite{liu2020meddg}。这类工作聚焦于文本输入，缺少多模态图像证据，也不包含"无法确诊即拒答"的安全机制。\textbf{本文的区别}在于：AquaDiag 是多模态（图像+主诉）的，且把合规红线、拒答判断与病灶识别作为一等评测维度。

\subsection{VLM 强化后训练与视觉遗忘}

可验证强化学习为多模态推理提供了一致接口。GRPO 通过组内相对得分估计优势，降低对独立价值模型的依赖~\cite{shao2024deepseekmath}；Reason-RFT 与 Patho-R1 分别在通用视觉推理与病理领域展示了 VLM 后训练的可行性~\cite{tan2025reasonrft,zhang2026pathor1}。TVC、CoFFT 等工作处理长推理中的视觉遗忘，通过在推理过程中重新注入视觉条件~\cite{sun2025tvc,zhang2025cofft}。\textbf{本文的区别}在于：我们关注的是"何时该追问、何时该拒答"的\textbf{行为对齐}，而非推理长度与效率；且 ProDiag 以可审计规则而非训练为主，直接证明规则结构即可修复行为缺陷。

\subsection{本文定位}

与只处理文本 RAG、只优化单一视觉基准、或只做澄清的工作不同，本文首次将\textbf{主动追问、合规红线、拒答判断与病灶识别}组合为低资源专业诊断的评测基准，并以"视觉显式化 + 规则表 + 红线过滤"的轻量方法证明：\textbf{可审计的领域规则结构能同时修复追问、合规、拒答三类行为}。我们强调可复现边界：基线统一多轮协议、指标全自动可审计、ProDiag 不训练、病灶本体受控。
